\documentclass[manuscript,screen 12pt, nonacm]{acmart}
\let\Bbbk\relax % Fix for amssymb clash 
\usepackage{vmlmacros}
\AtBeginDocument{%
  \providecommand\BibTeX{{%
    \normalfont B\kern-0.5em{\scshape i\kern-0.25em b}\kern-0.8em\TeX}}}
\usepackage{outlines}
\setlength{\headheight}{14.0pt}
\setlength{\footskip}{13.3pt}
\title{An Alternative to Pattern Matching, Inspired by Verse}
\author{Roger Burtonpatel}
\email{roger.burtonpatel@tufts.edu}
\affiliation{%
  \institution{Tufts University}
  \streetaddress{419 Boston Ave}
  \city{Medford}
  \state{Massachusetts}
  \country{USA}
  \postcode{02155}
}
\begin{document}
    
    \section{Equations}

    An intriguing alternative to pattern matching exists in~\it{equations}, from
    the Verse Calculus (\VC), a core calculus for the functional logic
    programming language~\it{Verse}~\citep{antoy2010functional,
    hanus2013functional, verse}. (For the remainder of this work, we~use “\VC”
    and “Verse” interchangeably.)

    Verse's equations and choice mechanism allow for much of the same expressive
    power as pattern matching with extensions. Rather than go into too much
    detail (the curious may examine the Verse paper, which is well worth a
    read), we support this claim with a familiar example: 


    But the VC semantics has a cost model that is hard to predict. Why? In~\VC,
    names (logical variables) are~\it{values}, and they can just as easily be
    the result of evaluating an expression as an integer or tuple. To bind these
    names,~\VC, like other functional logic languages, relies
    on~\it{unification} of its logical variables and \it{search} at runtime to
    meet a set of program constraints~\citep{antoy2010functional,
    hanus2013functional}. Combining unifying logical variables with search
    requires backtracking, which can lead to exponential runtime
    cost~\citep{hanus2013functional, wadler1985replace, clark1982introduction}. 

    % Todo: An example would be nice. Backtracking/evaluating backwards, then
    % done. 

    Pattern matching, by contrast, can be compiled to a~\it{decision tree}, a
    data structure that enforces linear runtime performance by guaranteeing no
    part of the scrutinee will be examined more than once~\citep{maranget}. A
    decision tree does not backtrack: once a program makes a decision based on
    the form of a value, it does not re-test it later with new information. 
    

    Figures: 

    \begin{figure}[ht] 
        \begin{minipage}[h]{0.54\linewidth}
          \verselst
          \begin{lstlisting}[numbers=none, basicstyle=\tiny, xleftmargin=.2em,
                            showstringspaces=false,
                            frame=single]
$\exists$ exclaimTall. exclaimTall = $\lambda$ sh. 
  one { 
      $\exists$ s. sh = $\langle$Square s$\rangle$; 
      s > 100.0; "Wow! That's a sizeable square!"
    | $\exists$ w h. sh = $\langle$Triangle, w, h$\rangle$; 
      h > 100.0; "Goodness! Towering triangle!"
    | $\exists$ b1 b2 h. sh = $\langle$Trapezoid, b1, b2, h$\rangle$;
      h > 100.0; "Zounds! Tremendous trapezoid!"
    | "Your shape is small." }
            \end{lstlisting}
                \Description{exclaimTall}
        \subcaption{\tt{exclaimTall} in~\VC }
            \label{fig:verseexclaimtall} 
        %   \vspace{4ex}
        \end{minipage}%%
        \begin{minipage}[h]{0.5\linewidth}
          \verselst
          \begin{lstlisting}[numbers=none, basicstyle=\tiny, xleftmargin=2em,
                        frame=single]
$\exists$ tripleLookup. tripleLookup = $\lambda$ rho x. 
    one { $\exists$ w. lookup rho x = $\langle$Just w$\rangle$; 
          $\exists$ y. lookup rho w = $\langle$Just y$\rangle$; 
          $\exists$ z. lookup rho y = $\langle$Just z$\rangle$;
          z
        | handleFailure x }
          \end{lstlisting}
                \Description{exclaimTall}
            \subcaption{\tt{tripleLookup} in~\VC }
            \label{fig:versetriplelookup} 
        \vspace{4ex}
        \end{minipage} 
        \begin{minipage}[h]{\linewidth}
          \verselst
          \begin{lstlisting}[numbers=none, basicstyle=\tiny, xleftmargin=9em, 
                            frame=single, showstringspaces=false]
$\exists$ game_of_token. game_of_token = $\lambda$ token. 
  $\exists$ f. one {
         token = one { $\langle$BattlePass, f$\rangle$ | $\langle$ChugJug, f$\rangle$ | $\langle$TomatoTown, f$\rangle$}; 
           $\langle$"Fortnite", f$\rangle$
       | token = one { $\langle$HuntersMark, f$\rangle$ | $\langle$SawCleaver, f$\rangle$}; 
           $\langle$"Bloodborne", 2 * f$\rangle$
       | token = one { $\langle$MoghLordOfBlood, f$\rangle$ | $\langle$PreatorRykard, f$\rangle$}; 
           $\langle$"Elden Ring", 3 * f$\rangle$
       |   $\langle$"Irrelevant", 0$\rangle$ }
          \end{lstlisting}
                \Description{game_of_token in~\VC}
        \subcaption{\tt{game\_of\_token} in~\VC }
            \label{fig:versegot} 
        \vspace{4ex}
        \end{minipage}%% 
    \caption{Code for the~\ref{pmandextensions} functions with equations looks
    similar, and it doesn't need extensions.}
    \label{fig:verseextfuncs}
      \end{figure}
        
    The code in Figure~\ref{fig:verseextfuncs} has all the Nice~Properties. This
    is promising for~\VC. If it rivals pattern matching with popular extensions
    in desirable properties, and~\VC does everything using only equations and
    choice, it seems like the language is an intriguing option for writing code! 

    \subsection{\VC has a challenging cost model}
    \label{vcbadcost}

    So what's the catch? In~\VC, names (logical variables) are~\it{values}, and
    they can just as easily be the result of evaluating an expression as an
    integer or tuple. To bind these names,~\VC, like other functional logic
    languages, relies on~\it{unification} of its logical variables and
    \it{search} at runtime to meet a set of program
    constraints~\citep{antoy2010functional, hanus2013functional}. Combining
    unifying logical variables with search requires backtracking, which can lead
    to exponential runtime cost~\citep{hanus2013functional, wadler1985replace,
    clark1982introduction}. 

    % Todo: An example would be nice. Backtracking/evaluating backwards, then
    % done. 

    Pattern matching, by contrast, can be compiled to a~\it{decision tree}, a
    data structure that enforces linear runtime performance by guaranteeing no
    part of the scrutinee will be examined more than once~\citep{maranget}. A
    decision tree does not backtrack: once a program makes a decision based on
    the form of a value, it does not re-test it later with new information. 
    

\end{document}
